\section{The original whole degree acts invertibly on every module}

\begin{theorem}[The whole-degree scalar automorphism]
\label{thm:kn-nadel-whole-degree-scalar-isomorphism}
\uses{lem:kn-nadel-whole-cover-degree,lem:kn-nadel-whole-cover-trace-degree}
\lean{EnvelopeProperty.Vanishing.NadelBoundedCanonicalDegreeScalar.originalDegree_nsmul_id_isIso,EnvelopeProperty.Vanishing.NadelBoundedCanonicalDegreeScalar.originalDegreeScalarIso,EnvelopeProperty.Vanishing.NadelBoundedCanonicalDegreeScalar.originalDegreeScalarIso_hom}
\leanok
Let $k$ be a field of characteristic zero and $f:X\to\operatorname{Spec}k$.
Let $F$ be the original divisor family, let $m_i>0$, and let $D$ be the
whole degree of the original joined root family, with every original
factor retained. For every $\mathcal O_X$-module $M$, multiplication by $D$
is an automorphism $d_D(M):M\simeq M$, whose underlying morphism is
$D\,\mathrm{id}_M$. The automorphism is constructed from $f$ and the
original positive degrees.
\end{theorem}
\begin{proof}
On every open $U\subseteq X$, the morphism $f$ gives the original
$k$-algebra structure on $\Gamma(U,\mathcal O_X)$. The positive whole
degree is nonzero in $k$ and hence is a unit there. Its image is therefore
a unit in $\Gamma(U,\mathcal O_X)$, including when $U$ is empty.

The native scalar-action ring homomorphism
\[
 \Gamma(U,\mathcal O_X)\longrightarrow
 \operatorname{End}\bigl(\Gamma(U,M)\bigr)
\]
sends this unit to the endomorphism $s\mapsto Ds$, which is consequently
invertible. The identity between this endomorphism and the $U$-component
of $D\,\mathrm{id}_M$ follows from additivity: for $D=0$ both are zero,
and increasing $D$ by one adds the identity on sections on both sides.
Thus every open-section component of $D\,\mathrm{id}_M$ is an
isomorphism. Isomorphisms of module sheaves are detected on these native
abelian-section components, so $D\,\mathrm{id}_M$ is an isomorphism.
Taking its associated isomorphism gives $d_D(M)$ with the stated
underlying morphism. No trace map or trace-normalization identity is
used.
\end{proof}
